%% ============================================================
%%  Referee Report
%%  Manuscript: "Significance Bias in Government Intervention
%%               Research: Evidence from Top Finance Journals"
%%  Journal: Journal of Finance (simulated review)
%% ============================================================

\documentclass[12pt]{article}
\usepackage[margin=1.25in]{geometry}
\usepackage{setspace}
\onehalfspacing
\usepackage{amsmath}
\usepackage{booktabs}
\usepackage{hyperref}
\hypersetup{colorlinks=true,linkcolor=black,citecolor=black,urlcolor=blue}
\usepackage{enumitem}
\usepackage{natbib}
\bibpunct{(}{)}{;}{a}{,}{,}
\usepackage{microtype}
\setlength{\parskip}{6pt}

\title{\textbf{Referee Report} \\[6pt]
       \large\textit{Significance Bias in Government Intervention Research:
       Evidence from Top Finance Journals} \\[6pt]
       \normalsize Submitted to the \textit{Journal of Finance}}

\author{Referee \#2 \\ Anonymous}

\date{May 2026}

\begin{document}
\maketitle
\thispagestyle{empty}

\noindent\rule{\linewidth}{0.4pt}

\section*{Summary Recommendation}

\textbf{Decision: Reject --- Encourage Resubmission}

The paper addresses an important and timely question: does the published
finance literature on government intervention systematically underrepresent
null findings?  If true, this would have significant implications for
regulatory policy, which relies heavily on the academic evidence base.  The
topic is well-motivated, the writing is generally clear, and the empirical
exercise is transparent.  However, the paper has several fundamental
identification problems that prevent me from recommending acceptance at this
stage.  The core issue is that the research design compares two populations
that are never directly observed at the same stage of the research pipeline,
which introduces confounds that the authors have not adequately addressed.
I describe these concerns below in order of severity.

\noindent\rule{\linewidth}{0.4pt}

\section{Major Concerns}

\subsection*{1. The Fundamental Identification Problem: Two Non-Overlapping Samples}

The paper's identification strategy rests on comparing publication rates
across direction codes within a combined sample of NBER working papers and
published journal articles.  The critical design choice is that only 8 of
183 published papers (4.4\%) are matched to an NBER working paper.  This
means the sample consists almost entirely of two non-overlapping populations:

\begin{itemize}
  \item 175 published journal articles with \emph{no} observed pre-publication
        counterpart.
  \item 113 NBER working papers with \emph{no} confirmed journal publication.
\end{itemize}

The paper's identifying assumption is that the NBER sample represents the
pre-publication distribution of findings and that the journal sample
represents the post-publication distribution.  But because these are
\emph{different papers} rather than the same papers tracked before and after
the editorial filter, the comparison is confounded.  The 113 unpublished
NBER papers and the 175 published journal papers may differ along many
dimensions other than direction code---including research quality, author
seniority, institutional affiliation, identification strategy, and data
quality.  If null-finding papers are of lower average quality (e.g., because
genuinely null results in this domain are harder to publish credibly due to
underpowering or weaker identification), the 36\% publication rate for null
papers reflects a quality gap, not a significance filter.  The paper does
not control for any of these alternative explanations.

The gold-standard design for this question would track the same papers from
pre-submission through editorial decision---as in \textit{Franco2014}'s
survey of NSF-funded studies, where grant records identify research regardless
of whether it was ultimately published.  The NBER working paper approach
approximates this, but the near-zero match rate means the approximation is
very rough.  The authors should either (a) dramatically improve the matching
rate, (b) add controls for observable quality proxies, or (c) substantially
qualify the causal interpretation.

\subsection*{2. NBER Selection Is Not a Random Draw from Research Output}

The paper acknowledges that ``NBER membership is itself selective'' but argues
the selection ``operates largely on researcher quality and institution rather
than on the direction of findings.''  This may be true in the cross-section,
but it creates a severe problem for the within-researcher inference the paper
requires.

Specifically: NBER-affiliated researchers are among the most productive and
well-connected academics in finance and economics.  Their null-finding papers
(which constitute 75.2\% of the NBER sub-sample) may be substantially more
publishable than null-finding papers from the broader population.  If this
is so, the 36\% publication rate for null-finding NBER papers represents an
\emph{upper bound} on the true publication probability for null results in the
field.  The paper's policy conclusion that ``the published record overstates
the share of interventions with detectable effects'' would then be
conservative---but the magnitude of the bias would also be smaller than
reported.  More problematically, there is no way to assess from the current
design whether the null papers in the NBER sample represent genuinely null
interventions or simply research programs that were ultimately abandoned
because early results were weak.

\subsection*{3. LLM Direction Coding: Validity and Potential Systematic Error}

The direction codes that drive all results are assigned by an LLM reading
the abstract of each paper.  Three concerns arise:

\textbf{(a) No systematic human validation.}  The paper mentions ``a
held-out set'' in the Related Literature section but provides no validation
statistics---no confusion matrix, no inter-rater agreement, no false-positive
or false-negative rates for direction coding.  Without these, readers cannot
assess whether the LLM's direction assignments are reliable.  A 10\% error
rate applied asymmetrically---for example, if the LLM more often miscodes
ambiguous positives as null than ambiguous negatives---could substantially
affect the estimated premiums.

\textbf{(b) Abstract language is itself affected by the publication process.}
Published abstracts are polished through the editorial process; NBER working
paper abstracts are often rougher and less definitive in their framing.  If
published papers are written to emphasize the significance and direction of
their results (which they are), the LLM may systematically assign stronger
direction codes to published papers and weaker or null codes to NBER papers.
This would mechanically generate the observed publication-rate gap even in the
absence of any true significance filter.

\textbf{(c) The definition of ``positive'' is normatively charged.}  A paper
finding that capital requirements reduce bank risk-taking is coded positive
(+1) because ``the intervention achieved its intended regulatory goal.''  But
many economists would code the same result as ``the intervention constrained
banks'' and ask whether the regulatory goal was desirable.  This coding
convention is not neutral, and it may map imperfectly onto what editors and
referees consider a ``positive'' finding.  The authors should test the
sensitivity of their results to alternative coding conventions.

\subsection*{4. Quality Differences Between Null and Directional Papers}

The most important alternative explanation for the null-paper disadvantage
is that null-finding papers in this domain are, on average, of lower empirical
quality.  Credibly estimating a null effect of a government policy requires
the same statistical power and identification rigor as estimating a positive
or negative effect---and in a domain where many interventions have genuinely
large effects, a finding of no effect is inherently harder to make credibly.
Editors and referees may rationally discount null results more heavily because
they are more likely to reflect underpowering than to reflect a true zero.

The paper does not control for any proxy of research quality: not
identification strategy (DiD vs.\ OLS), not number of citations, not author
quality, not data quality, not sample size, not the journal tier in which
NBER authors typically publish.  Until these confounds are addressed, the
significance-bias interpretation is not the only---or even the most
plausible---explanation for the observed pattern.

\subsection*{5. Inconsistency in Reporting the JF Results}

The paper makes a central claim that ``the finance publication process
exhibits strong significance bias.''  However, the journal-level results tell
a more nuanced story that is not adequately reconciled.  The JF coefficients
are labeled $1.022^{***}$ and $1.067^{***}$ in Table 4, but the text then
states they ``fall short of statistical significance at conventional levels''
and attributes this to low power.  There is a contradiction: triple-star
notation conventionally indicates $p < 0.01$, yet the text implies the JF
estimates are not statistically significant.  The authors should carefully
verify and reconcile the significance labels in Table 4.

More substantively, the fact that the largest journal in the field---and the
journal to which this paper is submitted---shows no significant significance
bias in this domain is a material finding that deserves more prominent
discussion rather than being explained away with a power argument.

\section{Moderate Concerns}

\subsection*{6. Sample Size and Statistical Power}

The analysis rests on 296 papers, of which 113 are NBER-only and 183 are
published.  The journal-level regressions have 151 (JF), 180 (RFS), and 191
(JFE) observations, but within each journal the published sub-sample is only
38, 67, and 78 papers, respectively.  The sub-period regressions have 97 and
154 observations.  These are small samples for probit models with multiple
parameters.  The authors should report that all results are verified with
linear probability models (which avoid separation issues in small samples)
and consider reporting bootstrap confidence intervals.

\subsection*{7. Interpretation of the 2000--2007 Sub-Period Omission}

The authors omit the 2000--2007 sub-period because it has only 45 papers
``insufficient for reliable probit estimation.''  But 45 papers is not an
unreasonably small sample for a bivariate probit, and the omission of this
sub-period---which covers the pre-crisis period when government intervention
was arguably less salient---could affect the time-trend interpretation.  At
minimum, the 2000--2007 results should be reported (with appropriate caveats)
rather than simply dropped.

\subsection*{8. The Corpus Description Is Internally Inconsistent}

The introduction states the authors screen ``a corpus of more than 40,000
documents,'' but the Data section reports a total corpus of 33,224 documents.
This discrepancy should be resolved.  (I note that the inconsistency appears
to stem from earlier draft language referring to the 1990--2024 period;
if the paper has been revised to cover only 2000--2024, all corpus references
must be updated consistently.)

\subsection*{9. Robustness Check R4 Is Misspecified}

Robustness check R4 restricts to the 169 high-confidence LLM classifications
and reports that the subsample is endogenously selected (88\% published,
versus 62\% in the full sample).  The authors correctly note that this renders
the coefficients uninformative and say they ``do not draw conclusions'' from
the exercise.  If that is the case, the check adds no information and should
either be replaced with a meaningful sensitivity analysis or omitted.

\section{Minor Comments}

\begin{enumerate}[leftmargin=*]

  \item \textbf{The NBER-to-journal match rate of 4.4\% is surprisingly low.}
  The paper attributes this to title revisions and non-NBER publication paths.
  The authors should quantify how many published papers they believe had NBER
  working paper precursors that were simply not matched (perhaps by checking
  known papers by hand), and assess whether the 175 unmatched published papers
  are a representative cross-section of the direction distribution.

  \item \textbf{No comparison to other domains.}  The paper would be
  strengthened considerably by a comparison sample---either from a different
  domain of finance (e.g., corporate governance without explicit policy
  intervention) or from another field (e.g., labor economics, where similar
  studies exist).  Without a comparison, it is hard to assess whether the
  documented 2.3$\times$ ratio is large or typical.

  \item \textbf{The discussion of mechanisms is speculative.}  The paper
  interprets the 2008--2015 negative-finding premium as ``elevated editorial
  demand for research documenting the unintended costs of the post-crisis
  regulatory expansion.''  This is a plausible narrative but is not tested.
  Editor letters, referee reports, or desk-rejection data could provide
  evidence for or against this mechanism.

  \item \textbf{No correction for multiple testing.}  The paper runs multiple
  probit specifications, journal subsets, and sub-period regressions.  Given
  the exploratory nature of some of these tests, Bonferroni or
  Benjamini--Hochberg corrections (or at minimum a discussion of multiple
  testing concerns) would strengthen the paper.

  \item \textbf{Terminology.}  The term ``significance bias'' is used
  throughout, but this term has a specific meaning in the statistics
  literature (bias in test statistics near the significance threshold, as in
  \textit{Brodeur2016}).  The phenomenon the authors document is more
  precisely described as a ``publication filter'' or ``null-result penalty.''
  Using established terminology or carefully defining the authors' usage
  would reduce confusion.

  \item \textbf{References.}  The paper cites \textit{Ash2023} for the use of
  LLMs in economics.  Given the rapid growth of this literature, the authors
  should also cite more recent work on LLM reliability and potential biases in
  text classification tasks, particularly for economics applications.

\end{enumerate}

\section*{Summary of Required Revisions}

For this paper to be publishable in a top journal, the authors must:

\begin{enumerate}
  \item Address the fundamental identification problem by substantially
        improving the NBER-to-journal match rate, adding controls for
        observable quality, or obtaining an alternative data source (e.g.,
        editor records, submission-level data from a journal) that allows
        tracking the same papers through the publication process.
  \item Provide systematic LLM validation statistics (confusion matrix,
        inter-rater reliability with human coders) for both the in-scope
        classification and the direction coding.
  \item Reconcile the JF significance labels in Table 4 with the text
        discussion.
  \item Report results for the 2000--2007 sub-period.
  \item Add controls for paper quality (author citations, identification
        strategy, sample size) or explicitly test whether the null-paper
        disadvantage survives quality controls.
  \item Replace or remove Robustness Check R4, which by the authors' own
        admission yields uninformative estimates.
\end{enumerate}

The topic is important and the paper is well-executed at the data-collection
and descriptive level.  The authors have built an impressive pipeline.
I encourage a revision that addresses the identification concerns, particularly
items 1--3 above.

\bigskip
\noindent\rule{\linewidth}{0.4pt}

\noindent\textit{Note: References in this report follow the style of the
manuscript under review.  ``Brodeur2016'' refers to Brodeur, L\'{e}, Sangnier,
and Zylberberg (2016); ``Franco2014'' refers to Franco, Malhotra, and
Simonovits (2014); ``Andrews2019'' refers to Andrews and Kasy (2019);
``Ash2023'' refers to Ash et al.\ (2023).}

\end{document}
